\documentclass[12pt,a4paper]{article}
\usepackage[utf8]{inputenc}
\usepackage[T1]{fontenc}
\usepackage{amsmath,amssymb}
\usepackage[margin=1in]{geometry}
\begin{document}
\title{Kuliah 08 - Metode Regresi Nonlinear}
\author{Yeni Herdiyeni}
\date{}
\maketitle

\maketitle
\tableofcontents
\newpage

\section{Kuliah 08 - Metode Regresi Nonlinear}

\textbf{Topik:} Regresi Polynomial, Fungsi Tangga (Step Function), dan Basis Function  

\textbf{Pengampu:} Yeni Herdiyeni, Program Studi Kecerdasan Buatan, IPB University  

\textbf{Tanggal:} 16 April 2026


\hrule


\subsection{1. Pendahuluan: Filosofi Regresi Nonlinear}

Dalam regresi linear sederhana, hubungan antara variabel independen $x$ dan variabel dependen $y$ diasumsikan berbentuk garis lurus:


\\[
\hat{y} = \beta_0 + \beta_1 x
\\]


Namun, banyak fenomena nyata - mulai dari pertanian, kesehatan, hingga industri - memiliki hubungan yang \textbf{tidak linear}. Misalnya, penambahan pupuk mungkin meningkatkan hasil panen pada awalnya, tetapi setelah titik optimal tertentu, penambahan lebih lanjut justru dapat menurunkan produksi karena keracunan nutrisi atau salinitas tanah.


\textbf{Inti dari regresi nonlinear} yang dibahas dalam kuliah ini bukanlah membuat model yang sulit diestimasi secara matematis. Sebaliknya, pendekatannya adalah:


\begin{quote}
\textbf{Membuat model tetap linear pada parameter, tetapi lebih fleksibel terhadap bentuk hubungan $x \rightarrow y$.}\end{quote}


Dengan kata lain, kita mengubah $x$ menjadi sekumpulan fitur baru $h_1(x), h_2(x), \dots, h_M(x)$, lalu melakukan regresi linear pada fitur-fitur hasil transformasi tersebut.


\subsubsection{Bentuk Umum Basis Function Regression}

\\[
\hat{y}(x) = \beta_0 + \sum_{m=1}^{M} \beta_m h_m(x)
\\]


Di mana:


Model ini tetap linear terhadap parameter $\beta$, sehingga estimasi dengan \textbf{Ordinary Least Squares (OLS)} masih dapat digunakan. Yang berubah adalah bahasa representasi dari $x$.


\subsubsection{Tiga Keluarga Penting}

\begin{tabular}{|l|l|l|l|}
\hline
No & Metode & Fungsi Basis & Karakteristik \\
\hline
1 & \textbf{Polynomial Regression} & $h_m(x) = x^m$ & Global, halus, mudah diinterpretasi \\
\hline
2 & \textbf{Step Function Regression} & $h_m(x) = \mathbb{I}(x \in A_m)$ & Lokal per interval, diskrit, mudah dibaca \\
\hline
3 & \textbf{Basis Function Regression} & Spline, RBF, Fourier, dll. & Lokal/fleksibel, kuat untuk pola kompleks \\
\hline
\end{tabular}


\begin{quote}
\textbf{Filosofi utama:} Hubungan data dapat dibuat fleksibel dengan memilih bahasa representasi yang tepat untuk $x$: global, lokal, halus, atau bertingkat.\end{quote}


\hrule


\subsection{2. Regresi Polynomial}

\subsubsection{2.1 Konsep dan Model Matematis}

Regresi polynomial memperluas regresi linear dengan menambahkan pangkat-pangkat dari $x$:


\\[
\hat{y}(x) = \beta_0 + \beta_1 x + \beta_2 x^2 + \beta_3 x^3 + \dots + \beta_d x^d
\\]


Di mana $d$ adalah derajat (degree) polynomial. Setiap pangkat $x^m$ memberikan kemampuan tambahan bagi kurva untuk "membengkok".


\textbf{Filosofi polynomial:} satu kurva global mencoba menjelaskan seluruh domain input.


\subsubsection{2.2 Transformasi Fitur}

Setiap data $x$ diubah menjadi vektor fitur baru:


\\[
x \rightarrow (x, x^2, x^3, x^4, x^5)
\\]


Dalam bentuk matriks desain:


\\[
X_{\text{poly}} =

\begin{bmatrix}

x_1 & x_1^2 & x_1^3 & x_1^4 & x_1^5 \\

x_2 & x_2^2 & x_2^3 & x_2^4 & x_2^5 \\

\vdots & \vdots & \vdots & \vdots & \vdots \\

x_n & x_n^2 & x_n^3 & x_n^4 & x_n^5

\end{bmatrix}
\\]


Kita menciptakan ruang fitur baru yang lebih kaya, tetapi estimasi parameter tetap menggunakan least squares.


\subsubsection{2.3 Estimasi Parameter}

\\[
\min_{\beta} \sum_{i=1}^{n} (y_i - \hat{y}(x_i))^2
\\]


Model tetap linear terhadap parameter, sehingga solusi closed-form OLS masih berlaku:


\\[
\hat{\beta} = (X_{\text{poly}}^T X_{\text{poly}})^{-1} X_{\text{poly}}^T y
\\]


\subsubsection{2.4 Use Case: Pertanian - Prediksi Hasil Panen}

Regresi polynomial cocok untuk memprediksi hasil panen berdasarkan dosis pupuk, intensitas penyiraman, atau umur tanaman. Hubungannya sering tidak linear:



Interpretasi:


\subsubsection{2.5 Implementasi Python}

\begin{lstlisting}[language=Python]
import numpy as np
from sklearn.pipeline import make_pipeline
from sklearn.preprocessing import PolynomialFeatures
from sklearn.linear_model import LinearRegression

# x perlu diubah menjadi matriks kolom
x = x.reshape(-1, 1)

model = make_pipeline(
    PolynomialFeatures(degree=5, include_bias=False),
    LinearRegression()
)

model.fit(x, y)
y_pred = model.predict(xx.reshape(-1, 1))
\end{lstlisting}

\begin{quote}
\textbf{Parameter kunci:} \texttt{degree} mengontrol fleksibilitas. Degree kecil menghasilkan bias tinggi (underfitting), sedangkan degree besar menghasilkan varians tinggi dan rawan overfitting.\end{quote}


\subsubsection{2.6 Kelebihan dan Kelemahan}

\textbf{Kelebihan:}


\textbf{Kelemahan:}


\hrule


\subsection{3. Regresi Fungsi Tangga (Step Function)}

\subsubsection{3.1 Konsep dan Model Matematis}

Step function membagi domain $x$ menjadi beberapa interval diskrit, dan setiap interval memiliki level prediksi yang relatif konstan:


\\[
\hat{y}(x) = \beta_0 + \sum_{m=1}^{M} \beta_m \mathbb{I}(c_{m-1} \leq x < c_m)
\\]


Di mana:


\textbf{Filosofi step function:} data dibagi menjadi beberapa interval, dan setiap interval memiliki level rata-rata sendiri.


\subsubsection{3.2 Use Case: Medis - Risiko Komplikasi Pasien}

Dalam praktik klinis, dokter sering menggunakan kategori ambang, misalnya kadar gula darah:


\begin{tabular}{|l|l|l|}
\hline
Kategori & Rentang & Interpretasi \\
\hline
Normal & $x < 100$ & Risiko rendah \\
\hline
Borderline & $100 \leq x < 126$ & Risiko sedang \\
\hline
Tinggi & $x \geq 126$ & Risiko tinggi \\
\hline
\end{tabular}


Model:


\\[
\hat{y} = \beta_0 + \beta_1 \mathbb{I}(x \in A_1) + \beta_2 \mathbb{I}(x \in A_2) + \beta_3 \mathbb{I}(x \in A_3)
\\]


Makna:


\subsubsection{3.3 Transformasi Fitur: Binning dan One-Hot Encoding}

Misalkan $x \in [0, 1]$ dibagi menjadi 8 interval:


\\[
A_1 = [0, 0.125), \quad A_2 = [0.125, 0.25), \quad \dots, \quad A_8 = [0.875, 1]
\\]


Setiap nilai $x$ diubah menjadi vektor indikator:


\\[
x \rightarrow (\mathbb{I}(x \in A_1), \mathbb{I}(x \in A_2), \dots, \mathbb{I}(x \in A_8))
\\]


Contoh: jika $x \in A_3$, maka representasinya adalah $(0, 0, 1, 0, 0, 0, 0, 0)$.


\subsubsection{3.4 Implementasi Python}

\begin{lstlisting}[language=Python]
from sklearn.pipeline import make_pipeline
from sklearn.preprocessing import KBinsDiscretizer
from sklearn.linear_model import LinearRegression

x = x.reshape(-1, 1)

model = make_pipeline(
    KBinsDiscretizer(n_bins=8, encode="onehot-dense", strategy="uniform"),
    LinearRegression()
)

model.fit(x, y)
y_pred = model.predict(xx.reshape(-1, 1))
\end{lstlisting}

\begin{quote}
\textbf{Parameter kunci:} \texttt{n_bins} menentukan resolusi model. Bin sedikit menghasilkan model terlalu kasar (high bias), sedangkan bin terlalu banyak membuat model sensitif terhadap noise (high variance).\end{quote}


\subsubsection{3.5 Kelebihan dan Kelemahan}

\textbf{Kelebihan:}


\textbf{Kelemahan:}


\hrule


\subsection{4. Regresi Basis Function}

\subsubsection{4.1 Konsep Umum}

Regresi basis function adalah kerangka umum yang mencakup polynomial dan step function sebagai kasus khusus. Bentuk umumnya:


\\[
\hat{y}(x) = \beta_0 + \sum_{m=1}^{M} \beta_m h_m(x)
\\]


Kita dapat memilih basis $h_m(x)$ sesuai dengan geometri dan karakteristik masalah:


\begin{tabular}{|l|l|l|}
\hline
Jenis Basis & Karakteristik & Cocok untuk \\
\hline
\textbf{Polynomial} & Global, halus & Pola mulus di seluruh domain \\
\hline
\textbf{Step/Indikator} & Lokal per interval & Sistem berbasis ambang \\
\hline
\textbf{Spline} & Halus dan lokal & Pola kompleks dengan kehalusan \\
\hline
\textbf{Gaussian RBF} & Lokal di sekitar pusat & Pola non-seragam, multi-modal \\
\hline
\textbf{Fourier} & Periodik & Data dengan pola musiman/siklis \\
\hline
\end{tabular}


\textbf{Filosofi basis function:} kita merancang basis sesuai geometri masalah. Basis lokal memberikan fleksibilitas tinggi tanpa harus menggoyang seluruh kurva.


\subsubsection{4.2 Use Case: Industri - Prediksi Kualitas Produk}

Dalam sistem industri, pola sering kompleks dan lokal:


Basis function cocok ketika pola data tidak cukup dijelaskan oleh satu kurva global sederhana.


\subsubsection{4.3 Radial Basis Function (RBF) Gaussian}

Salah satu basis lokal yang populer adalah Gaussian RBF:


\\[
h_j(x) = \exp\left(-\gamma (x - c_j)^2\right)
\\]


Di mana:


Setiap basis berbentuk lonceng (bell curve):


\subsubsection{4.4 Implementasi Python: RBF Custom Transformer}

\begin{lstlisting}[language=Python]
import numpy as np
from sklearn.pipeline import make_pipeline
from sklearn.linear_model import LinearRegression

class RBFFeatures:
    def __init__(self, centers, gamma=90):
        self.centers = np.asarray(centers)
        self.gamma = gamma

    def fit(self, X, y=None):
        return self

    def transform(self, X):
        X = np.asarray(X).reshape(-1, 1)
        return np.exp(-self.gamma * (X - self.centers) ** 2)


# Membuat 12 pusat basis di selang [0, 1]
centers = np.linspace(0, 1, 12).reshape(1, -1)

model = make_pipeline(
    RBFFeatures(centers=centers, gamma=90),
    LinearRegression()
)

model.fit(x, y)
y_pred = model.predict(xx.reshape(-1, 1))
\end{lstlisting}

\subsubsection{4.5 Representasi Matriks Fitur RBF}

\\[
\Phi =

\begin{bmatrix}

h_1(x_1) & h_2(x_1) & \dots & h_M(x_1) \\

h_1(x_2) & h_2(x_2) & \dots & h_M(x_2) \\

\vdots & \vdots & \ddots & \vdots \\

h_1(x_n) & h_2(x_n) & \dots & h_M(x_n)

\end{bmatrix}
\\]


Setiap kolom adalah satu Gaussian basis, dan setiap baris adalah representasi satu data dalam ruang basis.


\subsubsection{4.6 Model Matematis RBF Regression}

\\[
\hat{y}(x) = \beta_0 + \sum_{j=1}^{M} \beta_j \exp\left(-\gamma (x - c_j)^2\right)
\\]


Estimasi parameter tetap menggunakan least squares:


\\[
\min_{\beta} \sum_{i=1}^{n} (y_i - \hat{y}(x_i))^2
\\]


Model ini \textbf{linear terhadap parameter} tetapi \textbf{nonlinear terhadap variabel $x$}.


\hrule


\subsection{5. Perspektif Umum: Nonlinear sebagai Linear di Ruang Fitur Baru}

Semua metode yang dibahas memiliki satu kesamaan mendasar:


\begin{quote}
\textbf{Nonlinear regression ini sebenarnya adalah linear regression di ruang fitur baru.}\end{quote}


Perbedaannya terletak pada bagaimana kita merepresentasikan $x$:


\begin{tabular}{|l|l|l|}
\hline
Metode & Cara Berpikir & Representasi \\
\hline
\textbf{Polynomial} & Berpikir global & Satu cerita besar untuk semua data \\
\hline
\textbf{Step Function} & Berpikir diskrit & Memecah dunia menjadi kategori \\
\hline
\textbf{RBF/Basis Function} & Berpikir lokal & Banyak sudut pandang lokal yang saling overlap \\
\hline
\end{tabular}


\subsubsection{Pesan Penting}

Bukan modelnya yang berbeda, tetapi \textbf{cara kita merepresentasikan dan memahami data}.


\hrule


\subsection{6. Kapan Menggunakan Metode yang Mana?}

\begin{tabular}{|l|l|l|}
\hline
Polynomial & Step Function & Basis Function (RBF/Spline) \\
\hline
Pola halus global & Ingin segmentasi kasar & Pola kompleks dan lokal \\
\hline
Domain sempit & Perlu interpretasi interval & Ingin fleksibilitas terkontrol \\
\hline
Ingin model sederhana & Hubungan tidak perlu halus & Siap melakukan tuning basis \\
\hline
Efek berubah secara gradual & Sistem berbasis ambang & Data memiliki pola multi-modal \\
\hline
\end{tabular}


\hrule


\subsection{7. Penutup Filosofis}

Regresi nonlinear bukan sekadar menggambar kurva bengkok. Intinya adalah \textbf{mewakili realitas dengan basis yang sesuai}, lalu mengestimasi bobotnya secara sistematis.


Dengan memahami polynomial, step function, dan basis function, kita memiliki kerangka berpikir untuk memilih representasi data yang paling tepat sesuai konteks domain.


\hrule


\subsection{8. Ringkasan}

\begin{enumerate}
\item \textbf{Regresi fungsi tangga} membagi $x$ menjadi interval diskrit dengan prediksi konstan per interval.
\end{enumerate}

\begin{enumerate}
\item \textbf{Regresi basis function} menggunakan basis umum seperti RBF, spline, atau Fourier untuk fleksibilitas lebih tinggi.
\end{enumerate}

\begin{enumerate}
\item Semua metode tetap linear terhadap parameter, sehingga dapat diestimasi dengan OLS.
\end{enumerate}

\begin{enumerate}
\item Pemilihan metode bergantung pada karakteristik pola data: global, diskrit, atau lokal.
\end{enumerate}

\begin{enumerate}
\item Fleksibilitas harus diimbangi dengan waspada terhadap overfitting.
\end{enumerate}



\end{document}
